\iffalse
% \section{Дополнительные строковые алгоритмы: соглашения}
% Каждый листинг самостоятельный; общий префикс:
% \texttt{\#include <bits/stdc++.h>},
% \texttt{using namespace std; using ll=long long; using db=double;}.
% В листингах приведены только функции; все структуры локальны.
% Строки состоят из произвольных байтов, включая нулевой; лексикографический
% порядок задаётся значениями \texttt{unsigned char}, а не кодовыми точками Unicode.
% Индексы и длины имеют тип \texttt{int}, счётчики --- \texttt{ll}.
% Индексация с нуля, отрезки полуоткрытые. Пустые строки допустимы.
% Длины строк не превосходят $(\texttt{INT\_MAX}-3)/2$; память должна вмещать
% все создаваемые массивы и состояния, результаты счётчиков --- помещаться в \texttt{ll}.
% Все реализации написаны самостоятельно, а не скопированы из Algorithmica.
\fi
\section{Суффиксный массив удвоением и массив LCP}
Для строки байтов возвращает \texttt{suffixOrder[i]} --- начало $i$-го непустого суффикса в порядке \texttt{unsigned char}, и \texttt{lcp[i]} --- длину общего префикса суффиксов рангов $i$ и $i+1$.
Размеры ответов $n$ и $\max(0,n-1)$; пустая строка и нулевые байты допустимы, $n\le(\texttt{INT\_MAX}-3)/2$.
Время $O(n\log(n+1))$, память $O(n)$.
\begin{lstlisting}[name=extra_string_suffix_array]
pair<vector<int>, vector<int>> suffix_array_lcp(const string& text) {
    int n = int(text.size()), size = n + 1;
    vector<int> symbols(size), order(size), rank(size);
    for (int i = 0; i < n; ++i)
        symbols[i] = int(static_cast<unsigned char>(text[i])) + 1;
    vector<int> count(max(size, 257), 0);
    for (int symbol : symbols) ++count[symbol];
    for (int i = 1; i < 257; ++i) count[i] += count[i - 1];
    for (int i = size - 1; i >= 0; --i)
        order[--count[symbols[i]]] = i;
    int classes = 1;
    for (int i = 1; i < size; ++i) {
        if (symbols[order[i]] != symbols[order[i - 1]]) ++classes;
        rank[order[i]] = classes - 1;
    }
    vector<int> shifted(size), nextRank(size);
    for (int length = 1; length < size;) {
        for (int i = 0; i < size; ++i) {
            shifted[i] = order[i] - length;
            if (shifted[i] < 0) shifted[i] += size;
        }
        fill(count.begin(), count.begin() + classes, 0);
        for (int index : shifted) ++count[rank[index]];
        for (int i = 1; i < classes; ++i) count[i] += count[i - 1];
        for (int i = size - 1; i >= 0; --i)
            order[--count[rank[shifted[i]]]] = shifted[i];
        int nextClasses = 1;
        nextRank[order[0]] = 0;
        for (int i = 1; i < size; ++i) {
            int current = order[i], previous = order[i - 1];
            if (rank[current] != rank[previous] ||
                rank[(current + length) % size] !=
                rank[(previous + length) % size]) ++nextClasses;
            nextRank[current] = nextClasses - 1;
        }
        rank.swap(nextRank);
        classes = nextClasses;
        if (length >= size - length) break;
        length *= 2;
    }
    order.erase(order.begin());
    vector<int> position(n), lcp(max(0, n - 1));
    for (int i = 0; i < n; ++i) position[order[i]] = i;
    int matched = 0;
    for (int start = 0; start < n; ++start) {
        int place = position[start];
        if (place == n - 1) {
            matched = 0;
            continue;
        }
        int other = order[place + 1];
        while (start + matched < n && other + matched < n &&
               text[start + matched] == text[other + matched]) ++matched;
        lcp[place] = matched;
        if (matched > 0) --matched;
    }
    return {order, lcp};
}
\end{lstlisting}
% Сортирует непустые суффиксы строки. Возвращает пару \texttt{suffixOrder, lcp}:
% \texttt{suffixOrder[i]} --- начало $i$-го суффикса в лексикографическом порядке,
% \texttt{lcp[i]} --- длина общего префикса суффиксов с рангами $i$ и $i+1$.
% Размеры массивов равны $n$ и $\max(0,n-1)$. Используются удвоение с сортировкой
% подсчётом и алгоритм Касаи. Время $O(n\log(n+1))$, память $O(n)$.
% Служебный символ меньше любого байта и не входит в результат.
%
% RMQ по LCP даёт длину общего префикса произвольных суффиксов.
% Двоичный поиск по массиву находит вхождения образца; число различных
% непустых подстрок равно $n(n+1)/2-\sum lcp$, с вычислением произведения в \texttt{ll}.

\section{Суффиксный автомат: различные подстроки и вхождения}
Для строки \texttt{text} и образцов \texttt{patterns} возвращает число различных непустых подстрок и числа вхождений образцов в исходном порядке.
Строки состоят из произвольных байтов; перекрытия учитываются, пустой образец имеет $n+1$ вхождений, отсутствующий --- 0.
Длины не превосходят $(\texttt{INT\_MAX}-3)/2$, счётчики помещаются в \texttt{ll}; время $O((n+m)\log(\sigma+1)+q)$, память $O(n+q)$ при суммарной длине образцов $m$, их числе $q$ и $\sigma\le256$.
\begin{lstlisting}[name=extra_string_suffix_automaton]
pair<ll, vector<ll>> substring_counts(const string& text,
        const vector<string>& patterns) {
    struct State {
        int length = 0, link = -1;
        map<unsigned char, int> next;
        ll occurrences = 0;
    };
    int n = int(text.size()), last = 0;
    vector<State> states(1);
    states.reserve(2 * n + 1);
    for (unsigned char symbol : text) {
        int current = int(states.size());
        states.emplace_back();
        states[current].length = states[last].length + 1;
        states[current].occurrences = 1;
        int parent = last;
        while (parent != -1 && !states[parent].next.count(symbol)) {
            states[parent].next[symbol] = current;
            parent = states[parent].link;
        }
        if (parent == -1) states[current].link = 0;
        else {
            int target = states[parent].next.at(symbol);
            if (states[parent].length + 1 == states[target].length)
                states[current].link = target;
            else {
                int clone = int(states.size());
                State copy = states[target];
                copy.length = states[parent].length + 1;
                copy.occurrences = 0;
                states.push_back(move(copy));
                while (parent != -1) {
                    auto edge = states[parent].next.find(symbol);
                    if (edge == states[parent].next.end() ||
                        edge->second != target) break;
                    edge->second = clone;
                    parent = states[parent].link;
                }
                states[target].link = states[current].link = clone;
            }
        }
        last = current;
    }
    int size = int(states.size());
    vector<int> count(n + 1), order(size);
    for (const auto& state : states) ++count[state.length];
    for (int i = 1; i <= n; ++i) count[i] += count[i - 1];
    for (int i = size - 1; i >= 0; --i)
        order[--count[states[i].length]] = i;
    ll distinct = 0;
    for (int i = size - 1; i > 0; --i) {
        int vertex = order[i], parent = states[vertex].link;
        states[parent].occurrences += states[vertex].occurrences;
        distinct += states[vertex].length - states[parent].length;
    }
    vector<ll> answer;
    for (const string& pattern : patterns) {
        if (pattern.empty()) {
            answer.push_back(ll(n) + 1);
            continue;
        }
        int vertex = 0;
        for (unsigned char symbol : pattern) {
            auto edge = states[vertex].next.find(symbol);
            if (edge == states[vertex].next.end()) {
                vertex = -1;
                break;
            }
            vertex = edge->second;
        }
        answer.push_back(vertex == -1 ? 0 : states[vertex].occurrences);
    }
    return {distinct, answer};
}
\end{lstlisting}
% Одна функция строит автомат по \texttt{text} и возвращает пару:
% число различных непустых подстрок и массив чисел вхождений образцов
% из \texttt{patterns} в исходном порядке. Перекрытия учитываются;
% пустой образец имеет $n+1$ вхождений, отсутствующий --- ноль.
% Клоны не получают собственных вхождений; частоты распространяются по
% суффиксным ссылкам в порядке убывания длины.
% При $m$ --- суммарной длине образцов и $q$ --- их количестве время
% $O((n+m)\log(\sigma+1)+q)$, память $O(n+q)$, где $\sigma\le256$.
%
% Для малого фиксированного алфавита переходы можно хранить массивами.
% Сохранение позиции первого окончания позволяет восстанавливать подстроки;
% динамика по переходам позволяет искать $k$-ю лексикографическую подстроку.

\iffalse
\section{Палиндромное дерево}
Для строки произвольных байтов возвращает пары \texttt{length, occurrences} для различных непустых палиндромов в порядке первого появления.
Разные палиндромы одной длины представлены отдельно, перекрытия учитываются, служебные корни не возвращаются; пустая строка допустима.
Требуется $n\le(\texttt{INT\_MAX}-3)/2$, счётчики помещаются в \texttt{ll}; время $O(n\log(\sigma+1))$, память $O(n)$, $\sigma\le256$.
\begin{lstlisting}[name=extra_string_palindromic_tree]
vector<pair<int, ll>> palindrome_counts(const string& text) {
    struct Node {
        int length = 0, link = 0;
        map<unsigned char, int> next;
        ll occurrences = 0;
    };
    int n = int(text.size());
    vector<Node> nodes(2);
    nodes.reserve(n + 2);
    nodes[0].length = -1;
    int last = 1;
    auto extendable = [&](int vertex, int position) {
        while (true) {
            int left = position - 1 - nodes[vertex].length;
            if (left >= 0 && text[left] == text[position]) return vertex;
            vertex = nodes[vertex].link;
        }
    };
    for (int position = 0; position < n; ++position) {
        unsigned char symbol = static_cast<unsigned char>(text[position]);
        int parent = extendable(last, position);
        auto edge = nodes[parent].next.find(symbol);
        if (edge != nodes[parent].next.end()) last = edge->second;
        else {
            int current = int(nodes.size());
            nodes.emplace_back();
            nodes[current].length = nodes[parent].length + 2;
            nodes[parent].next[symbol] = current;
            if (nodes[current].length == 1) nodes[current].link = 1;
            else {
                int suffix = extendable(nodes[parent].link, position);
                nodes[current].link = nodes[suffix].next.at(symbol);
            }
            last = current;
        }
        ++nodes[last].occurrences;
    }
    for (int vertex = int(nodes.size()) - 1; vertex >= 2; --vertex)
        nodes[nodes[vertex].link].occurrences += nodes[vertex].occurrences;
    vector<pair<int, ll>> answer;
    for (int vertex = 2; vertex < int(nodes.size()); ++vertex)
        answer.emplace_back(nodes[vertex].length, nodes[vertex].occurrences);
    return answer;
}
\end{lstlisting}
% Возвращает по одной паре \texttt{length, occurrences} для каждого различного
% непустого палиндрома в порядке его первого появления при проходе слева направо.
% Разные палиндромы одной длины представлены отдельными элементами;
% число элементов --- число различных палиндромов. Вхождения могут перекрываться.
% Служебные корни длин $-1$ и $0$ не возвращаются. Время
% $O(n\log(\sigma+1))$, память $O(n)$, где $\sigma\le256$.
%
% Сохранение позиции первого окончания восстанавливает сами палиндромы.
% Серийные ссылки ускоряют динамику разбиения строки на минимальное число
% палиндромов; массивные переходы полезны при малом алфавите.

\fi
\section{Минимальный циклический сдвиг: алгоритм Бута}
Для строки возвращает наименьший индекс начала лексикографически минимального циклического сдвига; байты сравниваются как \texttt{unsigned char}.
Для пустой строки ответ 0; требуется $n\le(\texttt{INT\_MAX}-3)/2$.
Время $O(n)$, дополнительная память $O(1)$.
\begin{lstlisting}[name=extra_string_booth]
int minimal_cyclic_shift(const string& text) {
    int n = int(text.size());
    if (n == 0) return 0;
    int first = 0, second = 1, matched = 0;
    while (first < n && second < n && matched < n) {
        unsigned char a = static_cast<unsigned char>(text[(first + matched) % n]);
        unsigned char b = static_cast<unsigned char>(text[(second + matched) % n]);
        if (a == b) ++matched;
        else {
            if (a > b) {
                first += matched + 1;
                if (first <= second) first = second + 1;
            } else {
                second += matched + 1;
                if (second <= first) second = first + 1;
            }
            matched = 0;
        }
    }
    return min(first, second);
}
\end{lstlisting}
% Возвращает наименьший индекс начала лексикографически минимального
% циклического сдвига. Для пустой строки возвращает 0; при нескольких
% одинаковых минимальных сдвигах выбирается самый ранний.
% Время $O(n)$, дополнительная память $O(1)$: удвоенная строка не создаётся.
%
% Минимальный сдвиг задаёт каноническое представление ожерелья.
% Для максимального сдвига нужно обратить сравнение; сравнение также
% можно адаптировать к массивам целых чисел.

\iffalse
\section{Разложение Дюваля на слова Линдона}
Для строки байтов возвращает полуинтервалы множителей единственного разложения на лексикографически невозрастающую последовательность слов Линдона; сравнение по \texttt{unsigned char}.
Слово Линдона строго меньше каждого своего непустого собственного суффикса; одинаковые множители выдаются отдельно, пустой строке соответствует пустой ответ.
Требуется $n\le(\texttt{INT\_MAX}-3)/2$; время $O(n)$, дополнительная память $O(1)$ без ответа размера $O(n)$.
\begin{lstlisting}[name=extra_string_duval]
vector<pair<int, int>> lyndon_factorization(const string& text) {
    int n = int(text.size()), left = 0;
    vector<pair<int, int>> factors;
    while (left < n) {
        int right = left + 1, compared = left;
        while (right < n) {
            unsigned char a = static_cast<unsigned char>(text[compared]);
            unsigned char b = static_cast<unsigned char>(text[right]);
            if (a > b) break;
            if (a < b) compared = left;
            else ++compared;
            ++right;
        }
        int length = right - compared;
        while (left <= compared) {
            factors.emplace_back(left, left + length);
            left += length;
        }
    }
    return factors;
}
\end{lstlisting}
% Возвращает границы \texttt{[left, right)} единственного разложения строки
% на лексикографически невозрастающую последовательность слов Линдона.
% Слово Линдона строго меньше каждого своего непустого собственного суффикса.
% Одинаковые множители возвращаются отдельно. Пустой строке соответствует
% пустой ответ. Время $O(n)$, дополнительная память $O(1)$ без учёта ответа,
% ответ занимает $O(n)$.
%
% Разложение удвоенной строки позволяет найти минимальный циклический сдвиг.
% Границы можно обрабатывать сразу вместо хранения, если нужны только
% число множителей или агрегаты по ним.

\fi
